\subsection{An evaluated class-three Frobenius quotient of the cubic}

\begin{statement}{E.54}{Proved}{The genuine class-three action, its tangential transport, and every local lifting fiber at two}
Put $K=\mathbb Q(\zeta)$, $\zeta=\zeta_3$, $\lambda=1-\zeta$,
and $a=\zeta\lambda^{-3}$. Use the cubic $U_3$, its tangent
$s=z/x$ of value one at $O=[1:-1:0]$, and the actual map
\[
 h=\frac{x+y}{x+\zeta y}:U_3\longrightarrow
       \mathbb P^1\setminus\{0,1,\infty\}.
 \tag{E54.1}
\]
Let $N=\pi_1^{(3)}(U_{3,\overline K},\vec O)$ and let $F$ be
the free pro-three group on the standard positive peripheral
generators $x,y$ of the target. Choose the comparison path in
the punctured formal disc at zero, with compatible roots of $a$.
Then $h_*:N\twoheadrightarrow F$ is surjective. In particular
there is an actual arithmetic quotient
\[
 N/\gamma_4N\twoheadrightarrow C=F/\gamma_4F,\qquad
 \ker=\frac{h_*^{-1}(\gamma_4F)}{\gamma_4N}.
 \tag{E54.2}
\]
The maps from $N/N''$ and from
$N/\overline{\gamma_3(I_Q)}$ to $C$ factor through the
corresponding maps to $N/\gamma_4N$ in E.41.

Write
\[
 c=[x,y],\qquad d=[x,c],\qquad e=[y,c],
 \qquad [u,v]=uvu^{-1}v^{-1}.
 \tag{E54.3}
\]
Every element of $C$ has unique integral Hall coordinates
$x^A y^B c^n d^r e^s$, with all five exponents in $\mathbb Z_3$.
For $g\in G_K$, put $\chi=\chi_3(g)$ and $k=\kappa_a(g)$
using the chosen roots. The genuine based action on $C$ is
\[
 \begin{split}
 \alpha_{\chi,k}(x)&=x^\chi,\\
 \alpha_{\chi,k}(y)&=
       x^k c^{-b(\chi)}y^\chi c^{b(\chi)}x^{-k},
       \qquad b(\chi)=\frac{\chi^2-1}{24}\in\mathbb Z_3.
 \end{split}
 \tag{E54.4}
\]
Equivalently, all its Hall images are
\[
 \begin{aligned}
 x&\longmapsto x^\chi,\\
 y&\longmapsto
 y^\chi c^{k\chi}
 d^{\chi k(k-1)/2}
 e^{\,\chi b(\chi)-k\chi(\chi+1)/2},\\
 c&\longmapsto
 c^{\chi^2}
 d^{\,\chi^2(\chi-1)/2+k\chi^2}
 e^{\,\chi^2(\chi-1)/2},\\
 d&\longmapsto d^{\chi^3},\qquad
 e\longmapsto e^{\chi^3}.
 \end{aligned}
 \tag{E54.5}
\]
The composition law is
$\alpha_{\chi,k}\alpha_{\mu,l}=\alpha_{\chi\mu,k+\chi l}$.
For every compatible root choice its full image over $K$ is
the faithful image of the explicit group
\[
 \Gamma_a=
 \left\{(\chi,k)\in(1+3\mathbb Z_3)\ltimes\mathbb Z_3:
                 k\equiv\frac{\chi-1}{3}\pmod3\right\}.
 \tag{E54.6}
\]
Its outer image is the cyclotomic quotient, with the
factorization supplied by
$\alpha_{\chi,k}=\operatorname{Ad}(x^k)\alpha_{\chi,0}$.

Here is an exact local root choice and Frobenius evaluation.
At the inert prime of $K$ over two the residue field is
$\mathbb F_4$ and $\lambda$ has Teichmuller part $\zeta^2$.
Take
\[
 a_n=\zeta_{3^{n+1}}(a/\zeta)^{1/3^n},
 \tag{E54.7}
\]
where the second factor is the unique $3^n$th root in the
principal pro-two units. Arithmetic Frobenius $F_2$ then has
$\chi(F_2)=4$ and $\kappa_a(F_2)=1$ exactly. Thus
\[
 \begin{aligned}
 \alpha_{F_2}(x)&=x^4,&
 \alpha_{F_2}(y)&=y^4c^4e^{-15/2},\\
 \alpha_{F_2}(c)&=c^{16}d^{40}e^{24},&
 \alpha_{F_2}(d)&=d^{64},&
 \alpha_{F_2}(e)&=e^{64}.
 \end{aligned}
 \tag{E54.8}
\]
The alternate assembled roots
$a_n^{\rm old}=\zeta_{3^{n+1}}\lambda_n^{-3}$, with
$\lambda_n=\zeta_{3^{n+1}}^2
(\lambda/\zeta^2)^{1/3^n}$, give $\kappa_a(F_2)=-5$.
The actual transport between these two actions is
\[
 \frac{a_n}{a_n^{\rm old}}=\zeta_{3^n}^2,\qquad
 k_{\rm new}=k_{\rm old}+2(\chi-1),\qquad
 \alpha_{\rm new}
 =\operatorname{Ad}(x^{-2})\alpha_{\rm old}
                  \operatorname{Ad}(x^2).
 \tag{E54.9}
\]

Put $C_2=C/\gamma_3C$. All local $C_2$-valued cocycles at
$K_2$ are described by their values
\[
 \nu=x^a y^b c^n,\qquad z=x^A y^B c^N
 \tag{E54.10}
\]
on a positive tame inertia generator and arithmetic Frobenius,
respectively. The exact condition on the six coordinates is
\[
 L=4Ab-4Ba+6ab+4b+12n=0.
 \tag{E54.11}
\]
Use the same words as Hall lifts to $C$, with zero $d,e$
coordinates, and define
\[
 \begin{split}
 R={}&2A^2b-4ABa+16Aab+2Ab+16An
          -8Ba^2+2Ba-4Na\\
       &\hspace{12mm}+17a^2b+13ab+54an+40n,\\
 S={}&\bigl(16Ab^2-4Ab-4B^2a+4Ba+8Bb+32Bn-8Nb\\
       &\hspace{12mm}+14ab^2-6ab+16b^2+108bn+b+48n\bigr)/2 .
 \end{split}
 \tag{E54.12}
\]
The image of the map on local section classes
$H^1(K_2,C)\to H^1(K_2,C_2)$ is exactly the set of classes
whose representatives satisfy
\[
                 R,S\in3\mathbb Z_3.
 \tag{E54.13}
\]
For each such representative, one lift is specified by
\[
       t\longmapsto \nu d^{-R/60}e^{-S/60},\qquad
       F_2\longmapsto z.
 \tag{E54.14}
\]
Every nonempty geometric-conjugacy fiber has exactly $81$
elements. Its actual coordinates are the central changes
$z\mapsto zd^u e^v$, with $u,v\in\{0,\ldots,8\}$; its
translation group is $(\mathbb Z/9)^2$.

For a small explicit family take inertia $y$ and Frobenius
$x^{-1}c^N$ in $C_2$. These satisfy (E54.11), and their
class-three error is
\[
                 e^{\,5/2-4N}.
 \tag{E54.15}
\]
The members with $N\equiv1\pmod3$ lift. The member $N=0$
has empty fiber, whereas the member $N=1$ has the exact lift
with inertia $y e^{1/40}$ and Frobenius $x^{-1}c$.
They have the same abelian coordinates.

All the obstruction maps above are attached to (E54.2).
More precisely, for the local maps
$N/\gamma_4N\to C$ and $N/\gamma_3N\to C_2$, the diagram
of section sets commutes, and the central map is
$\operatorname{Lie}_3(N^{\rm ab})\to
 \operatorname{Lie}_3(F^{\rm ab})=\mathbb Z_3(3)^2$.
This gives an evaluated component of the intrinsic obstruction
in E.44. The displayed local $C_2$ classes are classes of this
specified quotient; their preimages in the full intrinsic
section set are governed by the kernel in (E54.2).
The further quotient of $C$ obtained by killing $\gamma_3C$
and the cubes of $x,y,c$ is exactly the order-$27$ quotient
of E.31 and E.35, with its stated based arithmetic action.
The manuscript passage is \S51.16, folios~733v--736v,
scans~172--178, equations~(371)--(379).
\end{statement}

\begin{proof}
\emph{The geometric quotient and its based action.}
The divisors computed directly in E.31 are
\[
 \operatorname{div}(h)=3O-3C_2,\qquad
 \operatorname{div}(1-h)=B_0+B_1+B_2-3C_2.
 \tag{E54.16}
\]
The two Kummer characters are independent modulo three.
Indeed the valuation at any $B_i$ first kills the coefficient
of $1-h$. A cube root of $h$ would have divisor $O-C_2$,
which would give a degree-one rational function on the
elliptic curve. The latter is impossible: a degree-one finite
map between smooth projective curves is an isomorphism, and
the genera are one and zero.
Thus the image of $N$ in $F/\Phi(F)$ is all of
$\mathbb F_3^2$. Its image in $F$ is closed, being the
continuous image of a compact group. A proper closed
subgroup of a finitely generated pro-three group is contained
in a maximal open subgroup of index three. Such a subgroup
cannot surject onto the Frattini quotient. This proves the
surjectivity and the exact kernel in (E54.2).
Surjectivity also identifies the image of each lower
central subgroup with the corresponding subgroup of $F$,
which proves the asserted central map.

The complete local expansion at $O$ is
$h=a s^3u(s)$, with $u(s)\in K[[s]]$ and $u(0)=1$.
Every compatible $3^n$th root of $u(s)$ with constant term
one belongs to $K[[s]]$, by formal binomial expansion.
The source tangent fixes the roots of $s$. In the punctured
target disc the only remaining Galois change is therefore
the power $q_g=x^{\kappa_a(g)}$. This statement precedes
passage to any nilpotent quotient: the geometric group of
the punctured formal disc is procyclic. Relative to the
standard target tangential action $\rho$, the induced
action is exactly
$\operatorname{Ad}(q_g)\rho_g$.
The arithmetic comparison is the continuous isomorphism
\[
 (C\rtimes_\alpha G_K)\longrightarrow
       (C\rtimes_\rho G_K),\qquad
       (w,g)\longmapsto(wq_g,g),
 \tag{E54.17}
\]
whose inverse is $(w,g)\mapsto(wq_g^{-1},g)$.
The cocycle identity
$q_{gh}=q_g\rho_g(q_h)$ proves the group law.

In the standard positive-path convention of D.23,
$\rho_g(x)=x^\chi$ and
$\rho_g(y)=f_g^{-1}y^\chi f_g$, with $f_g\in F'$.
Write $f_g=c^\beta$ modulo $\gamma_3F$.
Its genuine three-cycle relation, in the convention
of \cite[Theorem~2, equation~(5), and \S3]{A_Furusho2010Hexagon},
has $m=(\chi-1)/2$ and $xyz=1$.
In $F/\gamma_3F$ the three commutators
$[z,x],[y,z],[x,y]$ all equal $c$.
The remaining product $z^my^mx^m$ equals
$c^{-m(m+1)/2}$, by the class-two collection rule.
Consequently
$3\beta=m(m+1)/2$ and
$\beta=(\chi^2-1)/24$ with the positive sign.
Multiplication by three is injective on $\mathbb Z_3$,
and $\chi^2-1$ is divisible by three, so this is an
integral identity. The part of $f_g$ in $\gamma_3F$
is central in $C$ and drops out of its conjugation
action. This proves (E54.4), without choosing
unknown higher coefficients of the Galois conjugator.
The identity
$b(\chi\mu)=b(\chi)+\chi^2b(\mu)$ proves its composition
law, either on the generators or by (E54.5).

\emph{Integral Hall coordinates and the actual arithmetic image.}
The integral free nilpotent group of rank two and class three
has the Hall basis $x,y,c,d,e$. Collection gives a group
structure on $\mathbb Z_3^5$. If
$g=(a,b,n,r,s)$ and $g'=(A,B,N,R,S)$, its exact law is
\[
 \begin{split}
 gg'=\big(&a+A,\ b+B,\ n+N-Ab,\\
          &r+R+bA(A+1)/2-An,\\
          &s+S+Ab(b+1)/2-Bn+AbB\big).
 \end{split}
 \tag{E54.18}
\]
All denominators are units at three. The law is first
verified on integral powers using
$xy=cyx$, $xc=dcx$ and $yc=ecy$, with $d,e$ central,
and then extends continuously to $\mathbb Z_3$.
Uniqueness follows successively from the free integral
graded lattices of ranks $2,1,2$. Reduction modulo
$3^v$ in all five coordinates gives finite three-groups.
Their inverse limit is generated by $x,y$, and the
successive Hall coordinates show that its kernel from
$F$ is precisely $\gamma_4F$. This proves the asserted
integral model at the prime three; no division by three
in a Baker--Campbell--Hausdorff formula is used.
Applying (E54.18) to (E54.4) proves every entry of
(E54.5).

With roots assembled from $\lambda$ and $\zeta$,
\[
       k=\frac{\chi-1}{3}-3\kappa_\lambda.
 \tag{E54.19}
\]
Changing compatible roots adds $v(\chi-1)$ for some
$v\in\mathbb Z_3$, so it preserves the congruence in
(E54.6). The cyclotomic projection of $G_K$ is all of
$1+3\mathbb Z_3$. Over $L=K(\mu_{3^\infty})$, the
character $\kappa_\lambda$ has image $\mathbb Z_3$.
To prove this, $\lambda$ is not a cube in $L$:
$\lambda^2=-3\zeta$ and the existence of every
three-power root of $\zeta$ in $L$ would make a cube
root of three belong to $L$. But $L/\mathbb Q$ is
abelian, whereas $\mathbb Q(\sqrt[3]{3})/\mathbb Q$
is not Galois. The Kummer character is therefore
nonzero modulo three. Its closed image is a subgroup
of $\mathbb Z_3$ containing a unit and hence is
$\mathbb Z_3$. For a fixed cyclotomic value, (E54.19)
now gives precisely one coset of $3\mathbb Z_3$.
This proves equality with (E54.6).
Finally, the action on $C^{\rm ab}$ detects $\chi$.
For $\chi=1$ the action is $\operatorname{Ad}(x^k)$;
its effect on $y$ modulo $\gamma_3C$ detects $k$.
This proves faithfulness and the asserted outer image.

\emph{Frobenius and the change of roots.}
At the prime over two, reduction gives
$1-\zeta=\zeta^2$ in $\mathbb F_4$.
The map of taking a $3^n$th power is an automorphism
of the principal pro-two units. Thus the principal-unit
factors in (E54.7) and (E54.9) already lie in $K_2$
and are fixed by local Galois. Arithmetic Frobenius
acts on $\zeta_{3^{n+1}}$ by the fourth power.
It follows directly that
$\kappa_\lambda(F_2)=2$,
$\kappa_{a,{\rm old}}(F_2)=1-3\cdot2=-5$,
and $\kappa_{a,{\rm new}}(F_2)=1$.
Moreover
$\zeta_{3^{n+1}}^6=\zeta_{3^n}^2$, so the ratio of
the two root systems is exactly (E54.9).
Conjugating an action
$\operatorname{Ad}(x^k)\rho_g$ on the left by
$\operatorname{Ad}(x^{-2})$ and on the right by
$\operatorname{Ad}(x^2)$ changes $k$ to
$k+2(\chi-1)$, proving the transport with its sign.
Substitution of $(\chi,k)=(4,1)$ into (E54.5) gives
(E54.8). In general Hall coordinates that automorphism is
\[
 (a,b,n,r,s)\longmapsto
 \bigl(4a,4b,4b+16n,\ 40n+64r,\
       24n+64s-16\binom b2-15b/2\bigr).
 \tag{E54.20}
\]

\emph{All local lifts and their conjugacy fibers.}
The action on $C$ at $K_2$ is unramified: in the formula
(E54.4), both $\chi$ and the chosen unit Kummer character
are unramified. Its image is pro-three, as follows either
from (E54.6) or from its trivial action on the Frattini
quotient. A cocycle together with this action is a
homomorphism into a pro-three semidirect product.
Wild inertia, which is pro-two, therefore has zero image.
The actual maximal pro-three tame quotient is
\[
       \langle F_2,t\mid F_2tF_2^{-1}=t^4\rangle_{\rm pro\text{-}3}.
 \tag{E54.21}
\]
It is $\mathbb Z_3\rtimes\mathbb Z_3$ with action
$t\mapsto t^4$; the unit $4$ lies in $1+3\mathbb Z_3$,
so this action is continuous for the stated exponent
group. The tame generators and relation thus give both
the necessity and sufficiency of the cocycle relation
$z\alpha_{F_2}(\nu)z^{-1}=\nu^4$.

Use (E54.18) and (E54.20) for the two zero-central Hall
lifts in (E54.10). Direct collection gives
\[
 z\alpha_{F_2}(\nu)z^{-1}\nu^{-4}=(0,0,L,R,S)
 \tag{E54.22}
\]
in Hall coordinates, with exactly (E54.11)--(E54.12).
In particular the lower relation is $L=0$.
The central group is
$Z=\langle d,e\rangle=\mathbb Z_3(3)^2$,
and $F_2$ acts on it by $64$.
Changing the inertia lift by $d^r e^s$ changes
(E54.22) by $(64-4)(r,s)=60(r,s)$.
Changing the Frobenius lift by a central element
does not change that error.
Since $60\mathbb Z_3=3\mathbb Z_3$,
this proves the exact image condition and the
explicit correction (E54.14).

The differences between two successful lifts are
central one-cocycles with values in $Z$.
The tame relation forces their inertia value $v$
to satisfy $(64-4)v=0$, hence $v=0$.
Their Frobenius value is arbitrary in $Z$, and
a central conjugation changes it by $(1-64)w$.
Thus the central cohomology is exactly
$Z/63Z=(\mathbb Z/9)^2$.
There is no further quotient by a stabilizer of
the lower section. An element fixed by its
twisted action has, in its first nonzero lower
central graded coordinate, a vector fixed by
the scalar $4$ or $16$. Both scalar-minus-one
operators are injective on the corresponding
integral lattice. Its two coordinates therefore
vanish successively. This proves the fiber
cardinality and the representatives in the statement.
Substitution $(a,b,n,A,B,N)=(0,1,0,-1,0,N)$
in (E54.22) gives (E54.15), and
$60/40=3/2$ verifies the displayed positive lift.

Finally, all lower central quotients and their
obstruction cocycles are functorial for the
based arithmetic map (E54.2); applying $h_*$
to the actual relation defect is precisely
(E54.22). The order-$27$ projection has
$\chi\equiv1$ and $k\equiv(\chi-1)/3
=\kappa_3(\zeta)$ modulo three, so (E54.4)
reduces to the conjugation action in E.31.
This proves both comparison claims.

The exact programs
\texttt{a\_class3\_frobenius.py} and
\texttt{a\_class3\_local.py} verify the calculations
in an independent truncated noncommutative Magnus
algebra. They check the hexagon sign through degree
three, all Hall images, the affine composition law,
the actual inner root transport, the collection
law, and $639$ exact lower-section examples with
their central correction signs. The local example
and its $1/40$ lift are checked as literal group
identities. The polynomial derivations and the
integral arguments above establish the formulas
for all three-adic coordinates.
\end{proof}
